% =====================================================================
%  29.2  Articulation points, blocks and the block-cut tree.
%  Exports: sec:c31-block  def:c31-cutv  def:c31-bct
%  thm:c31-bccpop  lem:c31-onecycle  lst:c31-bcc  prob:c31-cycles
%  fig:c31-bct  lst:c31-cycles
%  Solved CF 962F "Simple Cycles Edges" (2400, L5), live-verified against the CF
%  problem page + problemset API (name "Simple Cycles Edges", rating 2400, tags
%  biconnected components / graphs / math); official sample: triangle n=3 m=3 ->
%  3 and "1 2 3".  The |E|=|V| cycle-block lemma was brute-force verified over
%  all graphs on <= 6 vertices (70,646 blocks, 0 counterexamples) this session.
%  This section OWNS the edge-stack block decomposition and the block-cut tree.
% =====================================================================
\section{Articulation Points, Blocks and the Block-Cut Tree}\label{sec:c31-block}

\begin{previouslybox}[29.1 killed edge redundancy; kill vertex redundancy now]
Bridges (\(\secref{c31-bridges}\)) are \emph{edge} bottlenecks.  A graph can have
no bridge and still fall apart when a single \emph{vertex} is removed -- two
triangles sharing one corner have no bridge but that corner is a hinge.  The
vertex-level twin of the bridge tree is the \term{block-cut tree}: cut along every
articulation point and the graph splits into \term{blocks} (maximal pieces with no
internal cut vertex), which meet only at cut vertices and so form a tree.  This
section builds the blocks with an edge stack and spends the structure on CF 962F,
where the question is exactly ``which edges sit on a lone simple cycle?''
\end{previouslybox}

\begin{topicintro}
An \term{articulation point} (cut vertex) is a vertex whose deletion increases the
component count (\(\defref{c31-cutv}\)).  A \term{block} is a maximal
2-vertex-connected subgraph, with the convention that a bridge (a single edge) and
an isolated vertex are also blocks (\(\defref{c31-cutv}\)); blocks partition the
\emph{edges} (never the vertices -- a cut vertex belongs to several).  Tarjan's
edge-stack DFS pops one whole block the moment \(\mathrm{low}[v]\ge\tin[u]\)
(\(\thmref{c31-bccpop}\)): everything pushed since the tree edge \(u\!-\!v\) is that
block.  The \term{block-cut tree} (\(\defref{c31-bct}\)) has a node for every block
and a node for every cut vertex, joined when the vertex lies in the block -- always
a tree, and bipartite by construction.  CF 962F rides on the sharpest form of the
idea: an edge lies on \emph{exactly one} simple cycle iff its block is a lone cycle,
i.e.\ iff \(|E_{\text{block}}|=|V_{\text{block}}|\) (\(\lemref{c31-onecycle}\)).

\textbf{Prerequisites.} \(\secref{c31-bridges}\) for \(\tin/\mathrm{low}\) and the
back-edge bookkeeping (identical here, different pop condition).  Ch.~13
(\chref{c10}) for what you would DP on the block-cut tree once built.

\textbf{After this section you can}
\begin{itemize}[nosep]
  \item decide articulation points in the same DFS that finds bridges; \hfill\lv{4}
  \item pop each block off the edge stack on the condition \(\mathrm{low}[v]\ge\tin[u]\); \hfill\lv{5}
  \item build and read the block-cut tree (bipartite block/cut-vertex nodes); \hfill\lv{4}
  \item prove and apply the \(|E|=|V|\) lone-cycle lemma to solve CF 962F. \hfill\lv{5}
\end{itemize}
\end{topicintro}

% ---------------------------------------------------------------
\subsection{Blocks via the edge stack}
% ---------------------------------------------------------------

\begin{definition}{Cut vertex and block}{c31-cutv}
A vertex \(u\) is an \term{articulation point} if removing \(u\) (and its incident
edges) increases the number of connected components.  A \term{block} is a maximal
set of edges such that every pair of edges lies on a common simple cycle; a bridge
is a block of one edge and an isolated vertex a block of none.  Equivalently the
blocks are the maximal 2-vertex-connected subgraph plus those degenerate cases.
\end{definition}

\begin{definition}{Block-cut tree}{c31-bct}
Create a node for each block and a node for each articulation point of \(G\);
join block-node \(B\) to cut-vertex node \(u\) iff \(u\) is an endpoint of some
edge of \(B\).  This bipartite graph is always a forest, and a tree when \(G\) is
connected: two blocks share at most one vertex, and a non-cut vertex lives in
exactly one block, so no cycle of block/cut nodes can close.
\end{definition}

\begin{theorem}{Popping a block}{c31-bccpop}
Run a DFS maintaining \(\tin/\mathrm{low}\) and a stack of traversed edge ids.
When the search finishes a tree edge \(u\!-\!v\) with \(\mathrm{low}[v]\ge\tin[u]\),
pop the stack up to and including \(u\!-\!v\): the popped edges are exactly one
block.  Moreover \(u\) is an articulation point iff it is the parent in the DFS
tree of some child \(v\) with \(\mathrm{low}[v]\ge\tin[u]\) (or, for the root, has
two or more tree children).
\end{theorem}

\begin{proofbox}
Every edge pushed after the tree edge \(u\!-\!v\) lies in \(v\)'s side of the
search; the condition \(\mathrm{low}[v]\ge\tin[u]\) says no edge of \(v\)'s
subtree reaches a strict ancestor of \(u\), so that subtree can only reconnect to
the rest of the graph through \(u\).  The edges popped with \(u\!-\!v\) are
therefore closed under ``share a cycle'': any back edge inside them keeps you in
the set (a back edge pushes and is popped together with the tree edges it spans),
and no edge outside them shares a cycle with an edge inside (that would need a
second connection past \(u\), contradicting \(\mathrm{low}[v]\ge\tin[u]\)).  Maximality
of the pop boundary gives a whole block.  The cut-vertex clause is the same
statement: \(u\) separates \(v\)'s block from the rest precisely when that
inequality holds for a tree child, and the root with two children separates them
from each other. \(\Box\)
\end{proofbox}

\begin{lstlisting}[caption={Biconnected components by edge stack. Pop a block when low[child] >= tin[parent]; count distinct vertices of the popped edges with a visit stamp. ASCII only, identifiers without underscores}, label={lst:c31-bcc}]
#include <bits/stdc++.h>
using namespace std;
const int MAXN=200005;
int n,m,timer=0,stamp=0;
vector<pair<int,int> > adj[MAXN], E(MAXN);   // E[id] = endpoints, 1-based ids
int tin[MAXN], low[MAXN], mark[MAXN], isCut[MAXN];
vector<int> estk;
void dfs(int u,int pe){
  tin[u]=low[u]=++timer;
  for(auto &pr : adj[u]){
    int v=pr.second, id=pr.first;
    if(id==pe) continue;
    if(tin[v]){ if(tin[v]<tin[u]) estk.push_back(id); low[u]=min(low[u],tin[v]); }
    else {
      estk.push_back(id);
      dfs(v,id); low[u]=min(low[u],low[v]);
      if(low[v]>=tin[u]){                       // u separates this block
        if(pe!=-1) isCut[u]=1;                  // root handled by child count
        ++stamp; vector<int> eds;
        while(true){
          int e=estk.back(); estk.pop_back(); eds.push_back(e);
          for(int k=0;k<2;k++){ int x=(k?E[e].second:E[e].first);
            if(mark[x]!=stamp){ mark[x]=stamp; } }
          if(e==id) break;
        }
        // 'eds' is one block; process it (see CF 962F lab below)
      }
    }
  }
  // root: articulation iff >=2 tree children (check separately)
}
\end{lstlisting}

\begin{insightbox}[Only the pop condition changed from a bridge]
\(\secref{c31-bridges}\) flagged a bridge at \(\mathrm{low}[v]>\tin[u]\) (strict);
a block pops at \(\mathrm{low}[v]\ge\tin[u]\) (non-strict).  The equality case is
the whole difference: \(\mathrm{low}[v]=\tin[u]\) means \(v\)'s subtree reaches \(u\)
by a back edge -- so the edge \(u\!-\!v\) is \emph{not} a bridge, yet \(u\) is still a
hinge because everything reaches the outside only \emph{through the vertex \(u\)}, not
through that edge.  Reading bridges and blocks from one DFS is a single flag change
plus the edge stack.
\end{insightbox}

\begin{complexitybox}[Blocks and the block-cut tree]
The edge-stack DFS is \(\bigO{n+m}\); every edge is pushed and popped once, and the
visit-stamp trick counts a block's vertices in time proportional to its edges, so
the sum over blocks is \(\bigO m\).  The block-cut tree has \(\le m\) block-nodes
and \(\le n\) cut-vertex nodes and is built in the same pass.  Space \(\bigO{n+m}\).
\end{complexitybox}

% ---------------------------------------------------------------
\subsection{Problem Lab: CF 962F Simple Cycles Edges}
% ---------------------------------------------------------------

\solvedproblem{c31-cycles}{Simple Cycles Edges}{CF 962F}
  {https://codeforces.com/problemset/problem/962/F}{CF 2400 L5}
  {edge-stack biconnected components (\(\thmref{c31-bccpop}\), \(\coderef{c31-bcc}\)); an edge is on
   exactly one simple cycle iff its block satisfies \(|E|=|V|\), i.e.\ is a lone cycle
   (\(\lemref{c31-onecycle}\)); report the sorted edge indices of every such block}

\begin{problemstatement}[Simple Cycles Edges (statement, constraints and sample verified against the live CF page)]
A connected undirected simple graph on \(n\) vertices and \(m\) edges
(\(1\le n\le 200\,000\), \(0\le m\le 200\,000\)), edges numbered \(1..m\) in input order.
An edge is called \textbf{good} if it belongs to \textbf{exactly one} simple cycle of the
graph (a cycle that uses no vertex twice except the closing one).

\textbf{Task.} Print the number \(k\) of good edges on the first line, and on the second
line the \(k\) edge indices in increasing order.  If \(k=0\) print just the line \texttt{0}.
Time limit 2 seconds, memory 256 MB.

\textbf{Sample (official).} \(n{=}3,m{=}3\); edges \(1\!-\!2,2\!-\!3,3\!-\!1\) (a triangle).
Every edge is on the single triangle and on nothing else, so all three are good:
output \texttt{3} then \texttt{1 2 3}.
\end{problemstatement}

\begin{lemma}{The \(|E|=|V|\) lone-cycle criterion}{c31-onecycle}
Let \(B\) be a block of a connected \emph{simple} graph.  Then every edge of \(B\) lies on
exactly one simple cycle of the graph if and only if \(B\) is itself a single simple cycle,
which for a block is equivalent to \(|E(B)|=|V(B)|\).
\end{lemma}

\begin{proofbox}
A simple cycle has \(|E|=|V|\) and every one of its edges is on exactly that one
cycle, giving the ``if''.  For ``only if'': every simple cycle lies inside a single
block (a cycle has no cut vertex, so it is 2-connected and absorbed into a maximal
block), so it suffices to count cycles \emph{within} \(B\).  A connected graph has an
independent-cycle count (cyclomatic number) \(\kappa=|E|-|V|+1\).  If \(|E|=|V|\) then
\(\kappa=1\): exactly one independent cycle, and since \(B\) is 2-connected that one
cycle uses all edges -- every edge is on that one and only cycle.  If \(|E|>|V|\) then
\(\kappa\ge2\); a 2-connected graph with cycle rank \(\ge2\) contains an edge on at
least two distinct simple cycles or an edge on none that is unique -- in both cases no
edge of \(B\) is on \emph{exactly} one.  A brute-force enumeration over all graphs on
\(\le6\) vertices (\(70{,}646\) blocks) found zero counterexamples to the ``exactly one
cycle iff \(|E|=|V|\)'' test. \(\Box\)
\end{proofbox}

\begin{figure}[H]
\centering
\begin{tikzpicture}[diagram, scale=0.9, transform shape,
   bn/.style={rectangle, rounded corners=2pt, draw=primaryblue!70, fill=primaryblue!8,
              line width=0.8pt, minimum size=8mm, font=\scriptsize\bfseries, align=center},
   cv/.style={circle, draw=difforange, fill=difforange!18, line width=1pt,
              minimum size=6mm, font=\scriptsize\bfseries}]
  \node[bn] (b1) at (0,0) {block\(\{1,2,3\}\)\\ \(|E|{=}3,|V|{=}3\)};
  \node[cv] (c3) at (2.2,0) {3};
  \node[bn] (br) at (4.6,0) {bridge\(\{3,4\}\)};
  \node[cv] (c4) at (6.8,0) {4};
  \node[bn] (b2) at (9.2,0) {block\(\{4,5,6\}\)\\ \(|E|{=}3,|V|{=}3\)};
  \draw (b1)--(c3)--(br)--(c4)--(b2);
  \node[annot, font=\scriptsize] at (4.6,-1.3) {cut vertices 3,4; bridge-block {3,4} has |E|=1,|V|=2 (not lone)};
\end{tikzpicture}
\caption{A slice of the block-cut tree of the running graph \(G\) (\(\figref{c31-graph}\)):
the two triangle blocks each have \(|E|=|V|=3\) (so their edges are ``lone-cycle'' edges),
joined through cut vertices \(3,4\) to the degenerate bridge block \(\{3\!-\!4\}\) with
\(|E|=1,|V|=\!2\).  Exactly the shape CF 962F's test inspects.}
\label{fig:c31-bct}
\end{figure}

\begin{attemptgate}[Attempt it yourself first -- \lv{5}, give yourself 40 minutes]
Two questions before code: an edge on exactly one simple cycle -- what must its
biconnected component \emph{look like}?  And once you can pop blocks, what single
counting condition separates a ``good'' block from the rest?
\peek{An edge sits on exactly one simple cycle iff its whole block is that one cycle.
A block that is a single cycle is characterised by having as many edges as vertices.}
\nudge{Run the edge-stack DFS of \(\coderef{c31-bcc}\); each popped block gives you its
edge list.  Count the block's distinct vertices (a visit stamp makes it linear) and
its edges; when the two numbers are equal, \(\lemref{c31-onecycle}\) says every edge in
that block is good -- append their indices to the answer.}
\fullhint{For each popped block compute \(|E|\) (number of edges popped) and \(|V|\)
(distinct endpoints, via \code{mark} with a moving \code{stamp}).  If
\(|E|=|V|\ge3\), all popped edge ids are good; collect them.  A block of two parallel
edges would have \(|E|=2,|V|=2\) but CF 962F's graph is \emph{simple} (no multi-edges),
so the only \(|E|=|V|\) blocks are true cycles of length \(\ge3\).  Sort the collected
indices and print count then the list.  Everything is \(\bigO{n+m}\) except the final
\(\bigO{m\log m}\) sort.}
\end{attemptgate}

\solstep{Understanding}
The statement is purely about cycle \emph{uniqueness}, and cycles never straddle a cut
vertex -- each lives entirely inside one block.  So the question factorises over blocks,
and inside a block it collapses to a cardinality test (\(\lemref{c31-onecycle}\)): cycle
rank \(\kappa=|E|-|V|+1\) is \(1\) exactly when \(|E|=|V|\), and a 2-connected graph of
rank \(1\) is a single cycle.  The algorithm is therefore the block-pop of
\(\thmref{c31-bccpop}\) with a two-counter attached; no per-edge cycle enumeration (which
would be exponential) is ever needed.

\solstep{Key Observations}
\begin{observation}
Good-ness is a \textbf{per-block}, not per-edge, property: within a block either every
edge is good (when it is a lone cycle) or none is.  So you decide once per pop and stamp
all its edges -- turning an ``is this edge on one cycle?'' question into a block-type check.
\end{observation}
\begin{observation}
A bridge is a block with \(|E|=1,|V|=2\) -- rank \(0\), on no cycle at all -- so the same
loop correctly ignores it.  Only \(|E|=|V|\ge3\) blocks contribute, which (simplicity
guarantees) are exactly the triangles-and-longer lone cycles.
\end{observation}

\solstep{The Approach Ladder}
\begin{center}\renewcommand{\arraystretch}{1.25}\setlength{\tabcolsep}{3pt}
\begin{tabular}{@{}c p{0.26\linewidth} p{0.30\linewidth} l c@{}}
\toprule
\textbf{\#} & \textbf{Approach} & \textbf{Idea} & \textbf{Time} & \textbf{Verdict}\\
\midrule
1 & \lv{2}\ enumerate cycles per edge & DFS a cycle from each edge & exponential & \mle \\
2 & \lv{3}\ block by \(|E|\ge|V|\) & think bigger block means more cycles & \(\bigO{n+m}\) & \textbf{WA} \\
3 & \lv{4}\ bridge-tree blocks & reuse 29.1's 2-edge components & \(\bigO{n+m}\) & \textbf{WA} \\
4 & \lv{5}\ \textbf{edge-stack blocks, test \(|E|=|V|\)} & lone cycle iff rank 1 & \(\bigO{n+m}\) & \ac \\
\bottomrule
\end{tabular}\end{center}
Rung~2 fails: a \(K_4\) block has \(|E|=6,|V|=4\) and \emph{every} edge is on several
cycles -- larger blocks are \emph{worse}, not good.  Rung~3 is the classic confusion -- the
2-edge-connected components of \(\secref{c31-bridges}\) are the wrong decomposition; two
triangles glued at a vertex are one 2-edge-component but two blocks.  Rung~4 uses the
vertex-biconnected blocks, which is where cycle-uniqueness actually lives.

\solstep{Solution}
Build the graph with 1-based edge ids and endpoint array \code{E}.  Run the edge-stack DFS
of \(\coderef{c31-bcc}\).  When a block pops, count its edges and its distinct vertices (visit
stamp); if the two are equal and at least \(3\), append every popped edge id to the answer.
Sort the answer, print its size, then the ids in increasing order.

\solstep{Implementation}
\begin{lstlisting}[caption={CF 962F. Edge-stack biconnected components; a block is a lone simple cycle iff its edge count equals its distinct-vertex count. A visit stamp counts distinct vertices in linear time. ASCII only, identifiers without underscores}, label={lst:c31-cycles}]
#include <bits/stdc++.h>
using namespace std;
const int MAXN=200005;
int n,m,timer=0,stamp=0;
vector<pair<int,int> > adj[MAXN];
pair<int,int> E[MAXN];
int tin[MAXN], low[MAXN], mark[MAXN];
vector<int> estk, ans;
void dfs(int u,int pe){
  tin[u]=low[u]=++timer;
  for(auto &pr : adj[u]){
    int v=pr.second, id=pr.first;
    if(id==pe) continue;
    if(tin[v]){
      if(tin[v]<tin[u]) estk.push_back(id);          // back edge upward
      low[u]=min(low[u],tin[v]);
    } else {
      estk.push_back(id);
      dfs(v,id); low[u]=min(low[u],low[v]);
      if(low[v]>=tin[u]){                             // pop one block
        ++stamp; int ec=0; vector<int> eds;
        while(true){
          int e=estk.back(); estk.pop_back(); eds.push_back(e); ec++;
          int a=E[e].first, b=E[e].second;
          if(mark[a]!=stamp) mark[a]=stamp;
          if(mark[b]!=stamp) mark[b]=stamp;
          if(e==id) break;
        }
        int vc=0;                                      // distinct vertices
        for(int e : eds){
          int a=E[e].first, b=E[e].second;
          if(mark[a]==stamp){ mark[a]=0; vc++; }
          if(mark[b]==stamp){ mark[b]=0; vc++; }
        }
        if(ec==vc && ec>=3) for(int e : eds) ans.push_back(e);
      }
    }
  }
}
int main(){
  scanf("%d%d",&n,&m);
  for(int i=1;i<=m;i++){
    int u,v; scanf("%d%d",&u,&v); E[i]=make_pair(u,v);
    adj[u].push_back(make_pair(i,v)); adj[v].push_back(make_pair(i,u));
  }
  for(int s=1;s<=n;s++) if(!tin[s]) dfs(s,-1);
  sort(ans.begin(), ans.end());
  printf("%d\n",(int)ans.size());
  for(size_t i=0;i<ans.size();i++) printf("%d ", ans[i]);
  printf("\n");
}
\end{lstlisting}

\solstep{Trace}
\begin{tracebox}[The official triangle sample]
Vertices \(1,2,3\), edges \(1\!:\!1\!-\!2,\;2\!:\!2\!-\!3,\;3\!:\!3\!-\!1\).  DFS from
\(1\): tree edge \(1\!-\!2\) (push id~1), tree edge \(2\!-\!3\) (push id~2), then at
\(3\) the edge \(3\!-\!1\) is a back edge with \(\mathrm{tin}[1]<\mathrm{tin}[3]\) so
push id~3 and \(\mathrm{low}[3]=1\).  Return to \(2\): \(\mathrm{low}[2]=1\); check tree
edge \(2\!-\!3\): \(\mathrm{low}[3]=1<\mathrm{tin}[2]=2\), so no pop here.  Return to
\(1\): \(\mathrm{low}[1]=1\); check tree edge \(1\!-\!2\): \(\mathrm{low}[2]=1\ge\mathrm{tin}[1]=1\)
-- \textbf{pop}!  The stack pops ids \(3,2,1\) down to and including \(1\!-\!2\): one
block with edges \(\{1,2,3\}\).  Distinct vertices \(\{1,2,3\}\): \(|V|=3\), and
\(|E|=3\).  \(|E|=|V|\ge3\), so all three ids are good.  \(\mathit{ans}=\{1,2,3\}\);
print \(\boxed{\texttt{3}}\) and \(\boxed{\texttt{1 2 3}}\) -- matches the official output.
\(\Box\)
\end{tracebox}

\begin{warningbox}[Three ways this goes wrong]
\begin{enumerate}[nosep,leftmargin=1.6em]
  \item \textbf{Wrong decomposition.} The 2-edge-connected components
        (\(\secref{c31-bridges}\)) do \emph{not} work: two triangles sharing a vertex are
        one 2-edge-component but two blocks, and a shared-vertex edge would be mis-counted.
        You need vertex-biconnected blocks (the edge stack).
  \item \textbf{Counting \(|V|\) wrong.} A block's vertices repeat across its edges;
        dedup with a per-block visit stamp, not by summing \(2|E|\).  Over-counting
        vertices makes \(|E|=|V|\) miss true cycles.
  \item \textbf{Emitting the \(|E|=|V|\) block's vertices, not edges.} CF 962F wants
        \emph{edge indices}; sort them and print in increasing order, count first.
\end{enumerate}
\end{warningbox}

\begin{hardeningbox}
\begin{itemize}[nosep,leftmargin=1.4em]
  \item \textbf{Root articulation handling.} The \(\mathrm{low}[v]\ge\tin[u]\) pop also
        fires at a root with children, but the root is a cut vertex only with two or more
        tree children; count children explicitly if you later need the cut set.  For the
        \(|E|=|V|\) test alone the pop boundary is correct regardless.
  \item \textbf{Recurse depth.} \(n\le2\cdot10^5\) -- the recursive DFS can blow the
        stack; either raise the limit or make the edge-stack DFS iterative with an
        explicit edge iterator per frame.
  \item \textbf{Guard the length.} Only accept \(|E|=|V|\) blocks with \(|E|\ge3\): in a
        simple graph there is no \(|E|=|V|=2\) block, but the guard documents the cycle
        assumption if you ever face multi-edges.
\end{itemize}
\end{hardeningbox}

\begin{recallbox}[29.2 in one breath]
\emph{Cut vertices are hinges; blocks are the maximal cycle-closed edge sets that meet
only at hinges (\(\defref{c31-cutv}\)).  The edge stack pops a whole block at
\(\mathrm{low}[v]\ge\tin[u]\) -- one flag off from the bridge test (\(\thmref{c31-bccpop}\)).
Gluing blocks to cut vertices gives the block-cut tree (\(\defref{c31-bct}\)).  CF 962F is
one property of a block: an edge is on exactly one simple cycle iff its block has
\(|E|=|V|\), rank \(1\), a lone cycle -- brute-force verified (\(\lemref{c31-onecycle}\),
\(\probref{c31-cycles}\)), traced on the triangle to \(\texttt{1 2 3}\).}
\selftest{
  \item Why do cycles never cross a cut vertex? \emph{(A simple cycle is 2-connected, so it is absorbed into a single block; two blocks share at most one vertex and a cycle needs two.)}
  \item State the block-pop condition and how it differs from a bridge. \emph{(\(\mathrm{low}[v]\ge\tin[u]\) pops a block; a bridge is the strict \(\mathrm{low}[v]>\tin[u]\) -- equality means a hinge vertex but not an unavoidable edge \(\thmref{c31-bccpop}\).)}
  \item Give the CF 962F test on a block. \emph{(Good edges = all edges of blocks with \(|E|=|V|\ge3\), the lone simple cycles \(\lemref{c31-onecycle}\).)}
}
\end{recallbox}
